\documentclass[11pt]{article}

\usepackage[a4paper,margin=2.5cm]{geometry}
\usepackage[ngerman]{babel}
\usepackage{amsmath,amssymb,amsthm,amsfonts}
\usepackage{mathtools}
\usepackage{hyperref}
\usepackage{enumitem}

\title{Priml\"ucken als Rotationen auf einem zyklischen Vier-Zustands-Raum}
\author{Thomas Hoffbauer}
\date{}

\newtheorem{theorem}{Satz}
\newtheorem{lemma}{Lemma}
\newtheorem{corollary}{Korollar}
\newtheorem{proposition}{Proposition}

\theoremstyle{definition}
\newtheorem{definition}{Definition}
\newtheorem{remark}{Bemerkung}

\begin{document}

\maketitle

\begin{abstract}
Wir betrachten Priml\"ucken $g_n=p_{n+1}-p_n$ zwischen aufeinanderfolgenden Primzahlen $>3$ modulo $12$. Die zul\"assigen Restklassen sind $\{1,5,7,11\}=(\mathbb{Z}/12\mathbb{Z})^\times$ mit $(\mathbb{Z}/12\mathbb{Z})^\times\cong C_2\times C_2$; verwendet wird jedoch nur eine zus\"atzlich gew\"ahlte zyklische Ordnung dieser vier Klassen. Der Hauptsatz zeigt, dass das zugeh\"orige Rotationsinkrement auf einem gerichteten Kreisgraphen $C_4$ ausschlie{\ss}lich von $g_n\bmod 12$ abh\"angt. Daraus folgt eine wohldefinierte Abbildung der sechs geraden L\"uckenklassen modulo $12$ auf vier Rotationszust\"ande. Die Konstruktion liefert damit eine exakte kombinatorische Kodierung von Priml\"ucken als diskrete Rotationsdynamik.
\end{abstract}

\section{Einleitung}

Priml\"ucken werden traditionell als arithmetische Gr\"o{\ss}en untersucht,
\[
g_n = p_{n+1}-p_n,
\]
wobei $p_n$ und $p_{n+1}$ aufeinanderfolgende Primzahlen sind.

Das Ziel dieser Notiz ist zu zeigen, dass dieselbe Information eine nat\"urliche geometrische Interpretation besitzt.
Anstatt Priml\"ucken als Abst\"ande auf der Zahlengeraden zu betrachten, interpretieren wir sie als Erzeuger von Rotationen auf einem zyklischen Vier-Zustands-System, das aus den zul\"assigen Restklassen modulo zw\"olf gewonnen wird.

Der mathematische Anspruch dieser Darstellung liegt in einer sauberen, exakten Umcodierung, nicht in einer neuen tiefen Struktur der Primzahlen. Insbesondere wird keine neue Gruppenstruktur eingef\"uhrt; vielmehr wird auf einer endlichen Vierermenge eine zus\"atzliche zyklische Ordnung gew\"ahlt, die die L\"uckenstatistik in eine dynamische Sprache \"ubersetzt.

\section{Der Restklassen-Zustandsraum}

Sei
\[
U_{12}
=
(\mathbb{Z}/12\mathbb{Z})^\times
=
\{1,5,7,11\}.
\]
Jede Primzahl $p>3$ geh\"ort genau zu einer dieser Restklassen.
Als Gruppe gilt $U_{12}\cong C_2\times C_2$, also ist $U_{12}$ nicht zyklisch.
Die Verwendung eines Zyklus bezieht sich hier nur auf eine zus\"atzliche Ordnung dieser vier Elemente.

\begin{definition}
Definiere die Zustandsabbildung
\[
\rho:U_{12}\rightarrow \mathbb{Z}_4
\]
durch
\[
\rho(1)=0,\qquad
\rho(5)=1,\qquad
\rho(7)=2,\qquad
\rho(11)=3.
\]
\end{definition}

Damit bestimmt jede Primzahl $p>3$ einen Zustand
\[
s(p)=\rho(p\bmod 12)\in\mathbb{Z}_4.
\]
Geometrisch identifiziert dies die zul\"assigen Restklassen mit dem gerichteten Kreisgraphen
\[
C_4:\qquad 0\rightarrow1\rightarrow2\rightarrow3\rightarrow0.
\]
\"Aquivalent dazu w\"ahlen wir auf den urspr\"unglichen Restklassen die zyklische Ordnung
\[
1\to5\to7\to11\to1.
\]

Aus historischen Gr\"unden verwenden wir gelegentlich auch die Bezeichnungen
\[
0=E,\qquad 1=A,\qquad 2=B,\qquad 3=C.
\]

\section{Priml\"ucken und Rotationsinkremente}

Seien $p_n<p_{n+1}$ aufeinanderfolgende Primzahlen und
\[
g_n=p_{n+1}-p_n.
\]
Dann gilt f\"ur $p_n>3$:
\[
g_n\equiv 0,2,4,6,8,10 \pmod{12}.
\]

\begin{definition}
Das zu $g_n$ geh\"orige Rotationsinkrement ist
\[
R_n=s(p_{n+1})-s(p_n)\pmod 4.
\]
\end{definition}

\begin{proposition}
Sei $\rho:U_{12}\to\mathbb{Z}_4$ die gew\"ahlte zyklische Nummerierung und
f\"ur $d\in\{0,2,4,6,8,10\}$ die Abbildung
\[
T_d:U_{12}\to U_{12},\qquad T_d(x)\equiv x+d\pmod{12}.
\]
Dann ist jedes $T_d$ eine Permutation von $U_{12}$, und es gibt genau ein $r(d)\in\mathbb{Z}_4$ mit
\[
\rho(T_d(x))\equiv \rho(x)+r(d)\pmod 4
\quad\text{f\"ur alle }x\in U_{12}.
\]
Explizit gilt
\[
r(0)=0,\quad r(2)=r(4)=1,\quad r(6)=2,\quad r(8)=r(10)=3.
\]
\end{proposition}

\begin{proof}
Da $d$ gerade ist, bildet $x\mapsto x+d\pmod{12}$ die vier ungeraden Klassen
$U_{12}=\{1,5,7,11\}$ wieder auf sich ab; auf einer endlichen Menge ist dies eine Permutation.
Die Werte von $r(d)$ ergeben sich durch direkte Auswertung von $T_d$ auf den vier Klassen:
\[
\begin{array}{c|c}
d \bmod 12 & \text{\"Ubergangsmuster } s\mapsto s+r(d)\ (\bmod 4)\\
\hline
0   & 0\to0,\ 1\to1,\ 2\to2,\ 3\to3\\
2   & 0\to1,\ 1\to2,\ 2\to3,\ 3\to0\\
4   & 0\to1,\ 1\to2,\ 2\to3,\ 3\to0\\
6   & 0\to2,\ 1\to3,\ 2\to0,\ 3\to1\\
8   & 0\to3,\ 1\to0,\ 2\to1,\ 3\to2\\
10  & 0\to3,\ 1\to0,\ 2\to1,\ 3\to2
\end{array}
\]
Damit ist die Formel f\"ur $r(d)$ gezeigt.
\end{proof}

\section{Darstellungssatz f\"ur Rotationen}

\begin{theorem}
Seien $p_n<p_{n+1}$ aufeinanderfolgende Primzahlen gr\"o{\ss}er als drei.
Dann h\"angt das Rotationsinkrement
\[
R_n=s(p_{n+1})-s(p_n)\pmod 4
\]
nur von der Restklasse der Priml\"ucke modulo zw\"olf ab.
Genauer gilt
\[
R_n=\Phi(g_n\bmod 12),
\]
wobei
\[
\Phi:\{0,2,4,6,8,10\}\rightarrow\mathbb{Z}_4
\]
durch
\[
\Phi(0)=0,\qquad
\Phi(2)=1,\qquad
\Phi(4)=1,\qquad
\Phi(6)=2,\qquad
\Phi(8)=3,\qquad
\Phi(10)=3.
\]
\end{theorem}

\begin{proof}
Setze $d=g_n\bmod 12\in\{0,2,4,6,8,10\}$.
Nach der vorigen Proposition ist die Wirkung von $x\mapsto x+d$ auf $U_{12}$
eine Permutation mit eindeutigem Inkrement $r(d)\in\mathbb{Z}_4$, so dass
\[
s(p_{n+1})-s(p_n)\equiv r(d)\pmod 4.
\]
Damit ist $R_n$ bereits eine Funktion nur von $d=g_n\bmod 12$, n\"amlich $R_n=r(d)$.
Mit der expliziten Tabelle aus der Proposition folgt:
\[
g_n\equiv0 \Rightarrow R_n=0,\qquad
g_n\equiv2,4 \Rightarrow R_n=1,\qquad
g_n\equiv6 \Rightarrow R_n=2,\qquad
g_n\equiv8,10 \Rightarrow R_n=3.
\]
Also h\"angt $R_n$ nur von $g_n \bmod 12$ ab.
\end{proof}

\section{Interpretation}

\begin{remark}
Der Begriff \emph{Rotation} wird hier graphentheoretisch verwendet:
Gemeint ist die orientierte Verschiebung auf dem gerichteten Kreisgraphen $C_4$,
nicht eine geometrische Rotation im euklidischen Raum.
\end{remark}

Der Satz liefert eine exakte Faktorisierung
\[
\text{Priml\"ucken}
\longrightarrow
\mathbb{Z}/12\mathbb{Z}
\longrightarrow
C_4.
\]
Damit erzeugt jede Priml\"ucke eine eindeutige Rotation auf dem Kreisgraphen.

Die Zuordnung lautet:
\[
\begin{array}{c|c|c}
g_n\bmod 12 & R_n & \text{Interpretation}\\
\hline
0 & 0 & \text{Keine Bewegung}\\
2,4 & 1 & \text{Vorw\"artsrotation}\\
6 & 2 & \text{Halbdrehung}\\
8,10 & 3 & \text{R\"uckw\"artsrotation}
\end{array}
\]

Folglich induziert die Primzahlfolge einen diskreten Rotationsfluss
\[
(s(p_n))_{n\ge 1}
\]
auf dem Kreisgraphen $C_4$.

\section{Primzahlvierlinge}

Betrachte einen Primzahlvierling
\[
(p,p+2,p+6,p+8).
\]
Sein Restklassenmuster modulo zw\"olf ist notwendig
\[
(5,7,11,1)
\]
oder
\[
(11,1,5,7).
\]

\begin{corollary}
Jeder Primzahlvierling induziert die Rotationsfolge
\[
(1,1,1),
\]
also drei aufeinanderfolgende Vorw\"artsrotationen auf $C_4$.
\end{corollary}

\begin{proof}
Die aufeinanderfolgenden L\"ucken sind
\[
2,\;4,\;2.
\]
Jede davon geh\"ort zur Klasse
\[
\{2,4\},
\]
die $R_n=1$ entspricht.
\end{proof}

\section{Dynamische Sicht}

Die Primzahlfolge
\[
(p_n)
\]
bestimmt daher eine Trajektorie
\[
(s(p_n))
\]
auf dem endlichen Zustandsraum $C_4$.
Priml\"ucken wirken als Bewegungserzeuger, w\"ahrend Primzahlkonstellationen als koh\"arente Rotationsmuster erscheinen.

Dies ergibt eine geometrische Neuformulierung der Priml\"uckenstatistik:
\[
\boxed{
\text{Priml\"ucken modulo }12
\Longleftrightarrow
\text{Rotationen auf }C_4.
}
\]

\section{Ausblick}

Ein naheliegender n\"achster Schritt ist eine nichttriviale dynamische Aussage \"uber die induzierte Folge $(R_n)$, etwa Asymmetrien zwischen Vorw\"arts- und R\"uckw\"artsinkrementen oder Grenzwertaussagen zu H\"aufigkeitsverh\"altnissen.

Eine pr\"azise testbare Anschlussaussage ist die folgende Vermutung:
\[
M_N
:=
\frac{\#\{1\le n\le N:\,R_n=1\}}
{\#\{1\le n\le N:\,R_n=3\}}.
\]
\emph{Vermutung.} Es existiert ein Wert $L\neq 1$ mit
\[
\lim_{N\to\infty} M_N = L.
\]
Diese Formulierung ist direkt numerisch \"uberpr\"ufbar und w\"urde die vorliegende Kodierung zu einer inhaltlichen dynamischen Aussage \"uber Priml\"ucken erweitern.

\section{Schluss}

Wir haben gezeigt, dass die zul\"assigen Restklassen der Primzahlen modulo zw\"olf eine nat\"urliche Darstellung als Vier-Zustands-Raum mit gew\"ahlter zyklischer Ordnung besitzen.
Unter dieser Darstellung induziert jede Priml\"ucke eine diskrete Rotation, deren Wert nur von der Restklasse der L\"ucke modulo zw\"olf abh\"angt.
Die resultierende Entsprechung ist eine exakte kombinatorische Umcodierung von Priml\"uckenklassen in ein diskretes Rotationssystem auf einem endlichen Kreisgraphen.

\end{document}